This work presents an experimental network intrusion detection laboratory focused on comparing Random Forest and Deep Neural Network (DNN) as supervised classifiers, evaluating the effects of data quality and dimensionality reduction via \textit{Mean Decrease in Impurity} (MDI) on the CICIDS2017 and LycoS-IDS2017 datasets. The initial hypothesis predicted that the DNN, due to its greater representational capacity, would achieve better generalization to attacks absent from training; however, this expectation was not confirmed. Given this limitation, an Autoencoder trained exclusively on benign traffic was investigated as an unsupervised alternative for unknown attack detection. Random Forest outperformed DNN in all evaluated scenarios, reaching a Macro F1 of 0.97 on LycoS-IDS2017 versus 0.87 for DNN, with inference time up to 5.18 times faster; the DNN showed difficulties with minority classes and greater sensitivity to data quality, gaining 0.20 Macro F1 points when replacing CICIDS2017 with LycoS-IDS2017, compared to 0.11 for Random Forest. MDI-based dimensionality reduction degraded all models by favoring majority-class features at the expense of discriminative characteristics for rare attacks. The Autoencoder achieved a Macro F1 of 0.89 and an attack Recall of 0.82 on LycoS-IDS2017, demonstrating partial viability for unlabeled scenarios, though without surpassing supervised classifiers on known attacks. It is concluded that Random Forest with curated data is the most suitable solution for production environments with known attacks, that dataset quality proves more decisive than architectural complexity, and that unknown anomaly detection remains an open challenge.

\noindent\textbf{Keywords:} Intrusion detection. \textit{Random Forest}. Deep neural networks. \textit{Autoencoder}. CICIDS2017.
